\documentclass[12pt,a4paper]{article}

\usepackage[T2A]{fontenc}
\usepackage[utf8]{inputenc}
\usepackage[russian,english]{babel}
\usepackage{amsmath, amssymb, amsfonts, mathtools}
\usepackage{booktabs}
\usepackage{siunitx}  
\usepackage{multirow}
\usepackage{longtable}

\title{Отчет о выполнении лабораторной работы по pthread}
\author{Антонов Артем Максимович 23КНТ3}


\begin{document}

\maketitle

\section{Множество Мондельброта}


\subsection{Математический алгоритм}

Множество Мандельброта --- это совокупность всех чисел $c \in \mathbb{C}$, для которых последовательность
\begin{equation}
  z_0 = 0, \qquad z_{n+1} = z_n^2 + c
  \label{eq:mandelbrot}
\end{equation}
является ограниченной для всех значений $n \in \mathbb{N}$. Иными словами, число $c$ принадлежит множеству Мандельброта тогда и только тогда, когда
\begin{equation}
  (\exists A \in \mathbb{R})\ (\forall n \in \mathbb{N})\ \ |z_n| < A.
  % \label{eq:membership}
\end{equation}

Представим $z^2$ как комплексное число, возведенное в квадрат:
\begin{align}
  &z = x + yi \\
  &z^2 = (x + yi)^2 = x^2 + 2xyi + y^2i^2 = x^2 + 2xyi - y^2 = (x^2 - y^2) + 2xyi
\end{align}
Прибавим к получившемуся выражению $c = a + bi$
\begin{align}
  &z^2 = (x^2 - y^2) + 2xyi + a + bi = (x^2 - y^2 + a) + (2xyb)i
\end{align}
В итоге получаем
\begin{align}
  z_{n+1} = (x^2 - y^2 + a) + (2xyb)i
\end{align}
Отсюда, 
\begin{align}
  &x_{n+1} = x_n^2 - y_n^2 + a \\
  &y_{n+1} = 2x_ny_nb
\end{align}

Теперь рассмотрим условие принадлежности числа $c$ к множеству Мондельброта:
\begin{align}
  |z_n| \le  2
\end{align}

Так как модуль комплексного числа это расстояние от точки до начала координат, можно $z$ представить так:
\begin{align}
  z_n = \sqrt {x_n^2 + y_n^2}
\end{align}
Далее подставим это в исходное выражение и возведемв квадрат
\begin{align}
  \sqrt {x_n^2 + y_n^2} < 2 \\
  x_n^2 + y_n^2 < 4
\end{align}

Итого, получаем, что при $x_n^2 + y_n^2 < 4$ точка $c$ принадлежит множеству Мондельброта.

Соответственно, при $x_n^2 + y_n^2 > 4$ - точка $c$ гарантированно не принадлежит множеству Мондельброта.


В итоге всех преобразований мы получили 3 основные формулы, которые будут использоваться в программе.

Для нахождения новый $x$ и $y$ на каждой итерации цикла проверки числа $c = a + bi$:
\begin{align}
  &x_{n+1} = x_n^2 - y_n^2 + a \\
  &y_{n+1} = 2x_ny_nb
\end{align}

И условие выхода из цикла проверки числа $c$:
\begin{align}
  x_n^2 + y_n^2 > 4
\end{align}
\subsection{Обзор кода}
Программа реализует следующий алгоритм нахождения точки множества Мондельброта:
\begin{enumerate}
\item С помощью генератора псевдослучайных чисел генерируем $a$ и $b$ и подгоняем их в диапазон точек.
\item Итерируемся $N$ раз, на каждой итерации генерируя $x$ и $y$ по формулам, которые были выведены в прошлой главе
\item Проверяем условия выхода из цикла, полученное в прошлой главе, если оно выполняется - ставим флаг и break
\item Проверяем флаг, если мы видим значение, означающее, что мы не вышли из цикла с помощью break (0), то записываем точку в глобальных массив точек.
\end{enumerate}

Этот алгоритм выполняется каждым потоком. Поток находит свое определенное количество точек, которое расчитывается по следующей формуле:

\begin{verbatim}
  int n_points = npoints / nthreads + (my_rank < npoints % nthreads ? 1 : 0);
\end{verbatim}

\subsection{Замеры производительности}
Замеры выполнялись на ubuntu 22.04 под WSL2. Процессор Intel Core i5-10400F - 6 ядер и 12 потоков.

Количество итераций $N$ в каждом запуске - 1000. Выполнялось по 15 запусков на 1, 2, 4, 8, 12, 16, 24 потоках. 
Количество точек, которое необходимо найти фиксированное - 1000000.

\begin{table}[htbp]
  \centering
  \caption{Результаты замеров}
  \label{tab:benchs}
  \begin{tabular}{cccrr}
    \toprule
    Run number & NPOINTS & Threads count & \multicolumn{1}{c}{AVG} & \multicolumn{1}{c}{MIN} \\
    \midrule
    15 & 1\,000\,000 & 1  & 6.3862 & 6.2632 \\
    15 & 1\,000\,000 & 2  & 3.2562 & 3.1483 \\
    15 & 1\,000\,000 & 4  & 1.7345 & 1.6990 \\
    15 & 1\,000\,000 & 8  & 0.9184 & 0.9019 \\
    15 & 1\,000\,000 & 12 & 0.7091 & 0.6618 \\
    15 & 1\,000\,000 & 16 & 0.7943 & 0.7655 \\
    15 & 1\,000\,000 & 24 & 0.6934 & 0.6736 \\
    \bottomrule
  \end{tabular}
\end{table}


\section{Своя реализация \texttt{rwlock}}

\subsection{Обзор кода}

В рамках лабораторной работы была реализована структура \texttt{my\_rwlock\_t} и основные  ее функции \texttt{rwlock}:
\begin{itemize}
  \item \texttt{int my\_rwlock\_init(my\_rwlock\_t* loc)}
  \item \texttt{int my\_rwlock\_destroy(my\_rwlock\_t* loc)}
  \item \texttt{int my\_rwlock\_rdlock(my\_rwlock\_t* loc)}
  \item \texttt{int my\_rwlock\_rwlock(my\_rwlock\_t* loc)}
  \item \texttt{int my\_rwlock\_unlock(my\_rwlock\_t* loc)}
\end{itemize}

\subsubsection{\texttt{my\_rwlock\_t}}
Структура \texttt{my\_rwlock\_t} содержит в себе следующие поля:
\begin{itemize}
  \item \texttt{pthread\_mutex\_t mutex} - мьютекс для внутренний операций. 
  \item \texttt{pthread\_cond\_t readers\_cond} - условная переменная для читателей.
  \item \texttt{pthread\_cond\_t writer\_cond} - условная переменная для писателей.
  \item \texttt{int readers\_counter} - счетчик активных читателей.
  \item \texttt{int wait\_read\_counter} - счетчик ожидающих читателей.
  \item \texttt{int wait\_write\_counter} - счетчик ожидающих писателей.
  \item \texttt{int writer\_active} - флаг, показывающий, есть ли активный писатель.
\end{itemize}
    

\subsubsection{\texttt{my\_rwlock\_init}}
Функция \texttt{my\_rwlock\_init} инициализирует поля структуры с помощью соответствующих стандартных функций, при ошибке инициализации возвращает код ошибки.

\subsubsection{\texttt{my\_rwlock\_destroy}}
Вызывает соответствующие \texttt{destroy} функции для объектов структур \texttt{pthread\_mutex\_t} и \texttt{pthread\_cond\_t}

\subsubsection{\texttt{my\_rwlock\_rdlock}}
Функция реализует механизм блокировки читателей. 
Для любых операций с внутренними переменными используются мьютекс. 
Соответственно, в начале функции мы блокируем мьютекс. 
Далее мы увеличиваем каунтер ожидающих читателей и запускаем цикл while с проверкой того, 
есть ли активный писать или есть ли ожидающие писатели, если условие выполняется - усыпляем поток, 
иначе уменьшаем количество ожидающих читателей и увеличиваем количество читателей и разблокируем мьютекс.

\subsubsection{\texttt{my\_rwlock\_wrlock}}
Функция реализует механизм блокировки писателей. Общая идея та же, мы увеличиваем каунтер и заходим в цикл, 
где условием является наличие активных писателей или наличие читателей. Если условие выполняется - засыпаем, иначе 
ставим флаг активного писателя и уменьшаем каунтер ожидающих писателей

\subsubsection{\texttt{my\_rwlock\_unlock}}
Общая функция разблокировки. Визуально поделена на 2 блока - в первом мы изменяем переменные, соответствующие
тому типа потока, который выходит, если это писать (смотрим по переменной \texttt{writer\_active}), то 
убираем активного писателя, иначе уменьшаем количество читателей. Следующий блок отвечает за передачу эстафеты, 
если у нас есть ожидающие писатели, и при этом нет активного писателя и количество читателей равно 0, то посылаем
сигнал спящему потоку-писателю, если ожидающих писателей нет - будим всех читателей.


\subsection{Замеры производительности}

Замеры выполнялись на ubuntu 22.04 под WSL2. Процессор Intel Core i5-10400F - 6 ядер и 12 потоков.
STDLIB - программа с использованием \texttt{rwlock} из pthreads, LAB - своя реализация.
Общие параметры для всех запусков: размер списка - 1000 элементов, количество операций - 80000
По 10 запусков на каждой тестовой конфигурации. Бралось минимальное и среднее время среди 10 запусков.

\begin{table}[htbp]
  \centering
  \caption{Результаты замеров}
  \label{tab:list-bench}
  \begin{tabular}{llrrrr}
    \toprule
    \multirow{2}{*}{\shortstack[l]{Search / insert /\\ delete, \%}}
      & \multirow{2}{*}{Threads}
      & \multicolumn{2}{c}{STDLIB} & \multicolumn{2}{c}{LAB} \\
    \cmidrule(lr){3-4} \cmidrule(lr){5-6}
     &  & min & average & min & average \\
    \midrule
    \multirow{6}{*}{99.9 / 0.05 / 0.05}
      & 1  & 0.0474 & 0.0493 & 0.0494 & 0.0527 \\
      & 2  & 0.0280 & 0.0289 & 0.0353 & 0.0363 \\
      & 4  & 0.0225 & 0.0233 & 0.0366 & 0.0382 \\
      & 8  & 0.0239 & 0.0257 & 0.0443 & 0.0457 \\
      & 16 & 0.0233 & 0.0245 & 0.0548 & 0.0565 \\
      & 24 & 0.0232 & 0.0240 & 0.0648 & 0.0661 \\
    \midrule
    \multirow{6}{*}{80 / 10 / 10}
      & 1  & 0.9064 & 1.0115 & 0.9578 & 1.0231 \\
      & 2  & 1.4652 & 1.4985 & 1.1762 & 1.2425 \\
      & 4  & 1.5506 & 1.5791 & 1.6709 & 1.6909 \\
      & 8  & 1.7337 & 1.7689 & 2.2691 & 2.2894 \\
      & 16 & 2.2558 & 2.3096 & 3.0469 & 3.0712 \\
      & 24 & 2.8347 & 2.8930 & 3.6787 & 3.7191 \\
    \midrule
    \multirow{6}{*}{50 / 25 / 25}
      & 1  & 3.5385 & 3.5817 & 3.5746 & 3.6058 \\
      & 2  & 4.9537 & 4.9961 & 4.1156 & 4.1923 \\
      & 4  & 5.0197 & 5.0684 & 4.7311 & 4.7634 \\
      & 8  & 5.2214 & 5.3020 & 5.2638 & 5.3256 \\
      & 16 & 6.9079 & 7.0078 & 7.1971 & 7.2405 \\
      & 24 & 8.3720 & 8.5332 & 8.8305 & 8.8971 \\
    \bottomrule
  \end{tabular}
\end{table}

\end{document}